% Generated from table-layout.md by md2tex.py, then reviewed against SKILL.md.

A column of short symbols next to a prose column: the prose column takes the spare width.

\needspace{202pt}

\begin{tabularx}{\textwidth}{|c|x{4.2cm}|L|c|}
    \hline
    \rowcolor{gray!15}
    \cc{Symbol} &
    \multicolumn{1}{y{4.2cm}|}{\textbf{Name}} &
    \multicolumn{1}{C|}{\textbf{Where it comes from}} &
    \multicolumn{1}{c|}{\textbf{Unit}} \\ \hline
\endhead

    $R$ &
    series resistance &
    the conductors are not perfect; current heats them &
    $\Omega/\text{m}$ \\ \hline

    $L$ &
    series inductance &
    current in the wires makes a magnetic field around them &
    H/m \\ \hline

    $G$ &
    shunt conductance &
    the insulation between the wires leaks a little current &
    S/m \\ \hline

    $C$ &
    shunt capacitance &
    the two wires are two conductors separated by insulation &
    F/m \\ \hline

\rowcolor{white}
\caption{The four line constants (per unit length)}
\label{tab:doc}
\end{tabularx}

A two-row table with short tall maths: kept on one page, and the short fraction column stays narrow.

\needspace{152pt}

\begin{tabularx}{\textwidth}{|c|x{3.4cm}|x{2.2cm}|L|}
    \hline
    \rowcolor{gray!15}
    \cc{Family} &
    \multicolumn{1}{y{3.4cm}|}{\textbf{Centre}} &
    \multicolumn{1}{y{2.2cm}|}{\textbf{Radius}} &
    \multicolumn{1}{C|}{\textbf{Special members}} \\ \hline
\endhead

    constant $r$ &
    \tallcell{$\left(\frac{r}{1+r}, 0\right)$ on the real axis} &
    \tallcell{$\frac{1}{1+r}$} &
    $r = 0$: the rim; $r = 1$: through the centre; $r \to \infty$: the point $(1, 0)$ \\ \hline

    constant $x$ &
    \tallcell{$\left(1, \frac{1}{x}\right)$ on the line $u = 1$} &
    \tallcell{$\frac{1}{\lvert x\rvert}$} &
    $x = 0$: the real axis; $x = \pm1$: centre $(1, \pm1)$, radius 1 \\ \hline

\rowcolor{white}
\caption{The two families of Smith-chart circles}
\label{tab:doc_2}
\end{tabularx}

A table longer than six rows may continue on the next page.

\begin{tabularx}{\textwidth}{|c|L|}
    \hline
    \rowcolor{gray!15}
    \cc{Step} &
    \multicolumn{1}{C|}{\textbf{Action}} \\ \hline
\endhead

    1 &
    Normalise the load impedance \\ \hline

    2 &
    Plot it on the chart \\ \hline

    3 &
    Draw the SWR circle through it \\ \hline

    4 &
    Move to the admittance point \\ \hline

    5 &
    Turn toward the generator to the g = 1 circle \\ \hline

    6 &
    Read the stub susceptance \\ \hline

    7 &
    Find the stub length from the short \\ \hline

\rowcolor{white}
\caption{Single-stub matching steps}
\label{tab:doc_3}
\end{tabularx}
